% Source PDF page 118; manual page 2-89.
\subsection{Range Insensitive Axis}\label{sec:range-insensitive-axis}

For a vehicle on a ballistic trajectory above the sensible atmosphere, an
instantaneous velocity increment $\Delta\vec V$ normally changes both time and
range to impact. If the increment is applied along a special direction called the
\emph{range-insensitive axis}, only the impact time changes; impact range remains
unchanged, as shown in \cref{fig:range-insensitive-axis}.

\begin{figure}[htbp]
  \centering
  \begin{tikzpicture}[>=Latex,x=1.0cm,y=1.0cm,line cap=round]
  \coordinate (v) at (0.6,2.65);
  \coordinate (p) at (8.5,0.0);
  \draw[line width=0.8pt] (v) .. controls (3.0,4.2) and (6.7,3.4) .. (p);
  \draw[line width=0.8pt] (v) .. controls (3.1,3.55) and (6.7,2.75) .. (p);
  \draw[line width=0.8pt] (v) .. controls (3.2,3.0) and (6.7,2.25) .. (p);
  \fill (v) circle (2pt);
  \fill (p) circle (2pt);
  \draw[->,line width=0.8pt] (v) -- ++(0,0.85);
  \draw[->,line width=0.8pt] (v) -- ++(0.65,-0.4);
  \node[left,align=right] at ($(v)+(-0.15,0.75)$)
    {$\Delta\vec V$ along range-\\insensitive axis};
  \node[below left] at (v) {Vehicle};
  \node[right,align=left] at (p) {Impact\\point};
  \node[above] at (5.9,3.3) {Ballistic trajectories};
  \draw[->] (6.35,3.05) -- (6.75,2.55);
\end{tikzpicture}

  \caption{Range-insensitive axis.}
  \label{fig:range-insensitive-axis}
\end{figure}

Two solutions exist. A velocity increment applied ``up'' the range-insensitive axis
increases time to impact; one applied ``down'' the axis decreases it.

The corresponding guidance rules align the body $x$ axis with one of these two
directions. TAOS first computes the initial-impact-point range and azimuth from the
current state with no velocity increment. It then uses a multidimensional
Newton--Raphson search to vary the direction of a
$\SI{10}{\foot\per\second}$ increment until the same impact point is recovered.
Aerodynamic forces are omitted from all impact-point calculations in this search.
The function \texttt{range\_insensitive\_axis} in the \texttt{derivs} module
returns the required pitch and yaw angles.
